% from ega2-1-fr.tex
\oldpage[II]{168}
Nous ne développerons pas davantage ce point de vue ici et laissons au
lecteur le soin de formuler de façon précise les définitions et résultats qui
correspondent aux indications précédentes.
\end{remark}

\subsection{Fermeture projective d'un fibré vectoriel}
\label{subsection:II.8.4-fr}

\begin{env}[8.4.1]
\label{II.8.4.1-fr}
Soient $Y$ un préschéma, $\mathcal{E}$ un $\mathcal{O}_Y$-Module
quasi-cohérent. Si on prend pour $\mathcal{S}$ la $\mathcal{O}_Y$-Algèbre
graduée $\mathbf{S}_{\mathcal{O}_Y}(\mathcal{E})$, la définition
(\hyperref[II.8.3.1.1-fr]{8.3.1.1}) montre que
$\widehat{\mathcal{S}}$ s'identifie à
$\mathbf{S}_{\mathcal{O}_Y}(\mathcal{E}\oplus\mathcal{O}_Y)$. Le cône affine
$\operatorname{Spec}(\mathcal{S})$ défini par $\mathcal{S}$ étant alors par
définition $\mathbf{V}(\mathcal{E})$, et $\operatorname{Proj}(\mathcal{S})$
étant par définition $\mathbf{P}(\mathcal{E})$, on voit donc que~:
\end{env}

\begin{proposition}[8.4.2]
\label{II.8.4.2-fr}
La fermeture projective d'un fibré vectoriel $\mathbf{V}(\mathcal{E})$ sur
$Y$ est canoniquement isomorphe à
$\mathbf{P}(\mathcal{E}\oplus\mathcal{O}_Y)$, et le lieu à l'infini de cette
dernière est canoniquement isomorphe à $\mathbf{P}(\mathcal{E})$.
\end{proposition}

\begin{remark}[8.4.3]
\label{II.8.4.3-fr}
Prenons par exemple $\mathcal{E}=\mathcal{O}_Y^r$ avec $r\geq2$~; alors les
cônes épointés $E$, $\widehat{E}$ définis par $\mathcal{S}$ ne sont ni affines
ni projectifs sur $Y$ si $Y\neq\varnothing$. La seconde assertion est
immédiate, car $\widehat{C}=\mathbf{P}(\mathcal{O}_Y^{r+1})$ est projectif sur
$Y$ et les espaces sous-jacents à $E$ et $\widehat{E}$ sont des ouverts non
fermés dans $\widehat{C}$, donc les immersions canoniques
$E\to\widehat{E}$ et $\widehat{E}\to\widehat{C}$ ne sont pas projectives
(\hyperref[II.5.5.3-fr]{5.5.3}), et on conclut par
(\hyperref[II.5.5.5-fr]{5.5.5, (v)}). D'autre part, supposant
$Y=\operatorname{Spec}(A)$ affine et par exemple $r=2$, on a
$C=\operatorname{Spec}(A[T_1,T_2])$ et $E$ est alors le préschéma induit par
$C$ sur l'ouvert $D(T_1)\cup D(T_2)$~; or nous avons vu que ce dernier n'est
pas affine (\hyperref[I.5.5.11-fr]{I, 5.5.11})~; 
\emph{a fortiori} $\widehat{E}$ ne peut être affine, puisque $E$ est l'ouvert
où ne s'annule pas la section $z$ au-dessus de $\widehat{E}$
(\hyperref[II.8.3.2-fr]{8.3.2}).

Toutefois~:
\end{remark}

\begin{proposition}[8.4.4]
\label{II.8.4.4-fr}
Si $\mathcal{L}$ est un $\mathcal{O}_Y$-Module inversible, on a des
isomorphismes canoniques pour les cônes épointés $E$ et $\widehat{E}$
correspondants à $C=\mathbf{V}(\mathcal{L})$,
\begin{equation}
\label{II.8.4.4.1-fr}
  \operatorname{Spec}\left(
    \bigoplus_{n\in\mathbf{Z}}\mathcal{L}^{\otimes n}
  \right)\xrightarrow{\sim}E
\tag{8.4.4.1}
\end{equation}
\begin{equation}
\label{II.8.4.4.2-fr}
  \mathbf{V}(\mathcal{L}^{-1})\xrightarrow{\sim}\widehat{E}.
\tag{8.4.4.2}
\end{equation}

En outre, il existe un isomorphisme canonique de la fermeture projective de
$\mathbf{V}(\mathcal{L})$ sur celle de $\mathbf{V}(\mathcal{L}^{-1})$,
transformant la section nulle (resp. le lieu à l'infini) de la première en le
lieu à l'infini (resp. la section nulle) de la seconde.
\end{proposition}

\begin{proof}
On a ici
$\mathcal{S}=\bigoplus_{n\geq0}\mathcal{L}^{\otimes n}$~; l'injection
canonique
\[
  \mathcal{S}\to\bigoplus_{n\in\mathbf{Z}}\mathcal{L}^{\otimes n}
\]
définit un morphisme dominant canonique
\begin{equation}
\label{II.8.4.4.3-fr}
  \operatorname{Spec}\left(
    \bigoplus_{n\in\mathbf{Z}}\mathcal{L}^{\otimes n}
  \right)
  \to\mathbf{V}(\mathcal{L})
  =\operatorname{Spec}\left(
    \bigoplus_{n\geq0}\mathcal{L}^{\otimes n}
  \right)
\tag{8.4.4.3}
\end{equation}
et il suffit de prouver que ce morphisme est un isomorphisme du schéma
$\operatorname{Spec}(\bigoplus_{n\in\mathbf{Z}}\mathcal{L}^{\otimes n})$
sur $E$. La question étant locale sur $Y$, on peut supposer que
$Y=\operatorname{Spec}(A)$ est affine

